\documentclass[11pt]{article}
\usepackage[margin=1in]{geometry}
\usepackage{amsmath,amssymb,amsthm,hyperref}
\newtheorem{theorem}{Theorem}
\newtheorem{lemma}{Lemma}
\newcommand{\leg}[2]{\left(\frac{#1}{#2}\right)}
\newcommand{\fr}[1]{\left\{#1\right\}}
\title{Quadratic remainder sums for two primes congruent to three modulo four}
\author{Principia Math In Progress}
\date{29 September 2026}
\begin{document}\maketitle
\begin{abstract}
We prove Chamberland--Dilcher Conjecture 7.2 for all positive exponents.
The proof expands the distribution of unit squares in quadratic characters
and partitions the original sum according to the two prime valuations.
This produces the conjectured class-number coefficients, including the
exceptional normalization at the prime three.
\end{abstract}

\section*{Task 1: the original conjecture}
Let $p,q$ be distinct primes congruent to $3$ modulo $4$, let $a,b\ge1$,
and set $n=p^a q^b$. Define
\[
 S_n=\frac1n\sum_{k=1}^{n-1}(k^2\bmod n)
     =\sum_{k=0}^{n-1}\fr{\frac{k^2}{n}},
 \qquad G_t(x)=\frac{x^t-1}{x-1}\quad(t\ge0).
\]
Here remainders are least nonnegative remainders and braces denote
fractional parts. Write
\[
 u=\lfloor a/2\rfloor,\quad v=\lfloor b/2\rfloor,\quad
 A=\lceil a/2\rceil,\quad B=\lceil b/2\rceil,
 \qquad H_r=h^*(-r).
\]
For a prime $r\equiv3\pmod4$, $H_r$ is the class number of
$\mathbb Q(\sqrt{-r})$ when $r>3$, while $H_3=1/3$.

\begin{theorem}\label{thm:main}
With this notation,
\begin{align*}
 S_n={}&\frac{n-p^u q^v}{2}\\
 &-G_A(p)\left(G_{v+1}(q)+\leg pq G_v(q)\right)H_p\\
 &-G_B(q)\left(G_{u+1}(p)+\leg qp G_u(p)\right)H_q.
\end{align*}
\end{theorem}
Since $A=\lfloor(a+1)/2\rfloor$ and $B=\lfloor(b+1)/2\rfloor$,
this is exactly Conjecture 7.2, Equation (7.2), in~\cite[p.~21]{CD}.
There is no restriction on the sizes or parities of the positive exponents.

\section*{Unit-square averages}
Let $\chi_r(x)=\leg xr$. We use the classical finite class-number formula
\begin{equation}\label{eq:class}
 \frac1r\sum_{x=1}^{r-1}x\chi_r(x)=-H_r.
\end{equation}
This is Dirichlet's formula as stated in~\cite[Theorem 4.1 and
Corollary 4.3]{CD}. At $r=3$ the sum is $(1-2)/3=-1/3$, which
explains the stated normalization.

For odd $M>1$ and $c$ coprime to $M$, put
\[
 U_M(c)=\sum_{\substack{x\bmod M\\(x,M)=1}}\fr{\frac{cx^2}{M}}.
\]
For $e,f\ge1$ the following formulas hold:
\begin{align}
 U_{p^e}(c)&=\frac{\varphi(p^e)}2-\chi_p(c)H_p,
 \label{eq:one}\\
 U_{p^e q^f}(1)&=\frac{\varphi(p^e q^f)}2
 -(1-\chi_p(q))H_p-(1-\chi_q(p))H_q.
 \label{eq:two}
\end{align}
We give the details so that their application to nonunit terms is explicit.

For an odd prime $r$, a unit residue $y$ modulo $r^e$ has
$1+\chi_r(y)$ square roots. Indeed, modulo $r$ it has zero or two
roots, and each root lifts uniquely at each step: solving
$(x+t r^j)^2\equiv y\pmod{r^{j+1}}$ is a linear equation for $t$
with nonzero coefficient $2x$ modulo $r$.
Consequently the number of roots of $cx^2=y\pmod{p^e}$ is
$1+\chi_p(c)\chi_p(y)$. Summing $y/p^e$ over unit residues $y$
gives~\eqref{eq:one}: the principal sum is $\varphi(p^e)/2$ by
pairing $y$ with $p^e-y$, and the character sum is $-H_p$ by
splitting into periods of length $p$ in~\eqref{eq:class}.

For $M=p^e q^f$, the Chinese remainder theorem gives the root count
\[
 (1+\chi_p(y))(1+\chi_q(y))
 \qquad ((y,M)=1).
\]
The principal term again sums to $\varphi(M)/2$. The product character
$\chi_p\chi_q$ is even and nonprincipal, so its weighted sum is zero:
pairing $y$ with $M-y$ reduces it to half its unweighted sum, which
vanishes. Finally, by periodicity and~\eqref{eq:class},
\[
 \frac1M\sum_{y=1}^{M}y\chi_p(y)=-H_p,
 \qquad
 \frac1M\sum_{\substack{1\le y\le M\\q\mid y}}y\chi_p(y)
 =-\chi_p(q)H_p.
\]
Subtracting excludes multiples of $q$; multiples of $p$ already contribute
zero. Interchanging $p$ and $q$ proves~\eqref{eq:two}.

\section*{Partition by valuations}
For $0\le k<n$, put
\[
 i=\min(v_p(k),a),\qquad j=\min(v_q(k),b),
\]
where the valuations of zero are taken to be infinite. In a fixed class
$(i,j)$, write $g=p^i q^j$, $m=n/g$, and $k=gx$. Then $x$ runs
once through the unit residues modulo $m$ (with the sole zero residue
when $m=1$).

Call the $p$ valuation active when $2i<a$, equivalently $0\le i<A$;
define active $q$ valuations analogously. If both are active, set
$M=p^{a-2i}q^{b-2j}$. Reduction from units modulo $m$ to units modulo
$M$ has constant fiber size $p^i q^j$, and $k^2/n=x^2/M$.
Thus this class contributes
\begin{equation}\label{eq:both}
 p^i q^j U_{p^{a-2i}q^{b-2j}}(1).
\end{equation}

If only the $p$ valuation is active, then $0\le i<A$ and $B\le j\le b$.
Now
\[
 \frac{k^2}{n}=\frac{q^{2j-b}x^2}{p^{a-2i}},
\]
and reduction to units modulo $p^{a-2i}$ has constant fiber size
$p^i\varphi(q^{b-j})$, with $\varphi(1)=1$. The contribution is
\begin{equation}\label{eq:p-only}
 p^i\varphi(q^{b-j})U_{p^{a-2i}}(q^{2j-b}).
\end{equation}
The formula with $p$ and $q$ interchanged treats the third case. If
neither valuation is active, then $n\mid k^2$ and the contribution is zero.
These four cases exhaust every residue $k$ without overlap.

\section*{Collecting the coefficients}
In each nonzero class, the principal term in~\eqref{eq:one} or
\eqref{eq:two}, multiplied by its fiber size, is half the number of
residues in that class. There are exactly $p^u q^v$ residues satisfying
$n\mid k^2$: they are the multiples of $p^Aq^B$ modulo $n$.
The total principal contribution is therefore $(n-p^u q^v)/2$.

Put $s=\chi_p(q)\in\{-1,1\}$. From~\eqref{eq:both}, the coefficient
of $-H_p$ in the classes with both valuations active is
\[
 G_A(p)(1-s)G_B(q).
\]
In~\eqref{eq:p-only}, the twisting character has constant value
$\chi_p(q^{2j-b})=s^b$. Since
\[
 \sum_{j=B}^{b}\varphi(q^{b-j})
 =\sum_{t=0}^{v}\varphi(q^t)=q^v,
\]
these classes contribute $G_A(p)s^b q^v$ to the same coefficient.
The classes with only $q$ active have no $H_p$ term. Hence its full
coefficient is
\[
 C_p=G_A(p)\big((1-s)G_B(q)+s^b q^v\big).
\]
Checking $b=2v$ and $b=2v+1$ separately, using
$G_{v+1}(q)=G_v(q)+q^v$, gives in both cases
\[
 (1-s)G_B(q)+s^b q^v=G_{v+1}(q)-sG_v(q).
\]
Quadratic reciprocity for $p\equiv q\equiv3\pmod4$ gives
$-s=\leg pq$. This yields the first class-number coefficient in
Theorem~\ref{thm:main}. Interchanging $p,a$ with $q,b$ gives the second.
Together with the principal contribution, this proves the theorem.

\section*{Tasks 2--5: follow-up questions}
\begin{enumerate}\setcounter{enumi}{1}
\item Derive a corresponding explicit valuation formula for three or more
distinct odd prime factors, tracking every odd quadratic character.
\item Extend the unit-square calculation to moduli divisible by arbitrary
powers of two, with explicit conductor and exceptional-unit corrections.
\item Prove uniform error bounds for $S_n-n/2$ as $n$ varies through
integers with a fixed number of distinct prime factors.
\item Determine exact formulas for higher moments of the normalized
remainders $\{k^2/n\}$ and identify the Dirichlet special values that occur.
\end{enumerate}

\paragraph{Verification.}
The accompanying integer/rational verifier evaluates the original sum,
enumerates reduced binary quadratic forms to check class numbers
independently, and checks the coefficient simplifications. Finite tests
corroborate the written all-exponent proof. Review is by the author-agent;
no independent peer review is claimed. Source and bounded-search limits
are recorded in the literature audit.

\begin{thebibliography}{9}
\bibitem{CD} M.~Chamberland and K.~Dilcher,
\emph{Sums of the floor function related to class numbers of imaginary
quadratic fields},
\href{https://arxiv.org/abs/2510.04387}{arXiv:2510.04387v1} (2025).
The complete current 23-page manuscript was read on 29 September 2026.
\bibitem{MathDB} MathDB,
\href{https://mathdb.com/p/370786}{problem 370786},
live record checked 29 September 2026.
\end{thebibliography}
\end{document}
